%%%% CAR-bench IJCAI-ECAI 2026 — Track 1 (Open) technical report.
%%%% Compile from a directory containing ijcai26.sty, named.bst, and track1_report.bib
%%%% (copy them from ../FormattingGuidelines-IJCAI-ECAI-26/, or build there).
%%%% 4-page content limit (competition rule). Keep \linenumbers ON for submission.

\documentclass{article}
\pdfpagewidth=8.5in
\pdfpageheight=11in

\usepackage{ijcai26}
\usepackage{times}
\usepackage{soul}
\usepackage{url}
\usepackage[hidelinks]{hyperref}
\usepackage[utf8]{inputenc}
\usepackage[small]{caption}
\usepackage{graphicx}
\usepackage{tikz}
\usetikzlibrary{arrows.meta}
\usepackage{amsmath}
\usepackage{booktabs}
\usepackage[switch]{lineno}

% Comment out for camera-ready.
\linenumbers

\urlstyle{same}

\pdfinfo{
/TemplateVersion (IJCAI.2026.0)
}

\title{Develop Small, Deploy Large: A Coroutine-Bridge Harness that\\
Transfers Reliability Across Model Scale on CAR-bench}

\author{
    Ivan Matveev
    \affiliations
    Proxima Ultra
    \emails
    i.a.matveev@gmail.com
}

\begin{document}

\maketitle

\begin{abstract}
CAR-bench measures whether tool-using agents stay reliable under real-world
uncertainty. We submit a coroutine-bridge harness whose only model action is to
emit a Python program that blocks and resumes in place across evaluator tool
exchanges, so deterministic policy lives as executable logic in the tool layer
rather than as prompt rules. Our Track~1 methodology treats model scale like a
train/test split: the harness is engineered and regression-tested against a
small, fast model (\texttt{gpt-oss-120b}) and then deployed unchanged on a
larger model (OpenAI \texttt{gpt-5.5}). Because the harness encodes transferable
task structure, the identical agent runs on the larger model with no
architecture change and resolves the residual cases the small model handled
unreliably. The design's efficiency is model-independent: two model calls carry
a full multi-turn task on both models. Since the larger model's calls are an
order of magnitude more expensive in latency, the bridge's call-minimization
matters more the larger the model. The single static prompt also stays
byte-identical across calls and across tasks, so most input tokens are served
from provider prefix caches --- a property shared by Cerebras (73\% cached
across our development runs) and the OpenAI route this track uses.
\end{abstract}

\section{Introduction}

CAR-bench evaluates the \emph{consistency} and \emph{limit-awareness} of LLM
agents across 58 interconnected tools, 19 domain policies, and 254 public tasks
in base, hallucination, and disambiguation categories~\cite{kirmayr2026carbench}.
The Open track ranks agents by hidden-set Pass\textsuperscript{3}, which credits
a task only when all three trials pass, and additionally rewards innovation in
harnessing and reliability design.

Our submission is built on one idea applied twice. The model's only action is to
emit an executable Python program~\cite{wang2024codeact}; a coroutine bridge
suspends that program when it calls a tool, emits the official Agent-to-Agent
(A2A) tool call, and resumes the same stack frame when results return --- with no
new model call. Because the action surface is code, deterministic CAR-bench
policies are compiled into the tool layer as ordinary logic rather than written
as prompt rules the model must remember.

The Track~1 contribution is a \emph{methodology}, stated as an analogy to a
train/test split over model scale. We ``train'' the harness --- hand-engineer it
and lock it with regression tests --- against a small, cheap, fast model where
iteration is inexpensive, and ``test'' it by deploying the identical agent on a
larger model. A change is accepted only if it improves the structure both models
can use, never if it patches a weakness specific to the small one. The result is
an agent whose reliability transfers upward across a large capability gap without
re-engineering.

\section{The Coroutine Bridge}

CAR-bench externalizes tool execution to the evaluator so it can score the tool
trajectory and inject environment perturbations. Every batch of tool calls the
agent emits returns as a fresh A2A message. A conventional next-action agent
therefore spends one model call per sequential tool step. The bridge removes this
tax: the model emits one program, and tool results are delivered back into the
blocked Python frame, which resumes without a model invocation
(Figure~\ref{fig:bridge}). One program drives many sequential and parallel tool
round-trips per model call. Independent reads are issued together through
\texttt{batch([...])}, which collapses independent round-trips into one parallel
A2A message. Every tool call still appears in the official trajectory, so the
technique is fully compliant.

\begin{figure}[t]
\centering
\resizebox{\linewidth}{!}{%
\begin{tikzpicture}[
  font=\small,
  msg/.style={-{Stealth[length=2mm]}, thick},
  ret/.style={-{Stealth[length=2mm]}, thick, dashed},
  call/.style={-{Stealth[length=2mm]}, very thick, blue!70!black},
  callret/.style={-{Stealth[length=2mm]}, very thick, dashed, blue!70!black},
  life/.style={dashed, gray!60},
  hdr/.style={draw, rounded corners, fill=gray!12, minimum height=6mm,
              inner sep=4pt, font=\small\bfseries},
  note/.style={draw, fill=yellow!22, rounded corners, font=\footnotesize,
               inner sep=2.5pt},
  lbl/.style={font=\footnotesize, above, midway},
  step/.style={font=\footnotesize\itshape, fill=blue!6},
]
\def\xE{0}\def\xW{4.2}\def\xM{8.4}\def\xP{13.2}
\node[hdr] (E) at (\xE,0) {Evaluator};
\node[hdr] (W) at (\xW,0) {Worker};
\node[hdr] (M) at (\xM,0) {Model};
\node[hdr] (P) at (\xP,0) {Python frame $+$ Bridge};
\foreach \x in {\xE,\xW,\xM,\xP} \draw[life] (\x,-0.45) -- (\x,-8.0);
\draw[msg] (\xE,-0.9) -- node[lbl]{user message (A2A)} (\xW,-0.9);
\node[step] at ({(\xW+\xM)/2},-1.35) {model call 1 of $\leq$5};
\draw[call]    (\xW,-1.9) -- node[lbl]{execute\_python} (\xM,-1.9);
\draw[callret] (\xM,-2.5) -- node[lbl]{Python program} (\xW,-2.5);
\draw[msg] (\xW,-3.1) -- node[lbl]{run program} (\xP,-3.1);
\draw[msg] (\xP,-3.9) -- node[lbl]{tool\_calls (A2A)} (\xE,-3.9);
\draw[ret] (\xE,-4.5) -- node[lbl]{tool\_results (A2A)} (\xP,-4.5);
\node[note, anchor=east] at (\xP-0.25,-5.0) {resume --- no model call};
\draw[msg] (\xP,-5.6) -- node[lbl]{dependent tool\_call (A2A)} (\xE,-5.6);
\draw[ret] (\xE,-6.2) -- node[lbl]{tool\_results (A2A)} (\xP,-6.2);
\node[note, anchor=east] at (\xP-0.25,-6.7) {resume --- no model call};
\draw[ret] (\xP,-7.3) -- node[lbl]{observation / respond()} (\xW,-7.3);
\end{tikzpicture}}
\caption{One user turn. Blue = model call; black solid = request, dashed =
return. A single \texttt{execute\_python} program drives several evaluator tool
round-trips; each returning tool result resumes the blocked Python frame with no
new model call. The loop is bounded at five sequential model calls (Track budget).}
\label{fig:bridge}
\end{figure}

\paragraph{Policy as code.}
Many CAR-bench policies are deterministic given grounded facts: whether opening a
window past 25\% needs an air-conditioning confirmation, whether high beams are
blocked while fog lights are on, which route default applies to a new segment.
The code-action surface lets us encode these as logic in the tool layer ---
guarded wrappers that fire whether the model calls the public name or the helper,
a confirmation state machine that stores the fully grounded pending action,
live tool-surface membership checks for missing capabilities, and result
normalization that gives stable field names. The runtime owns facts and strict
policy; the model owns interpretation of the request. A policy in the prompt can
be forgotten under attention pressure; a policy in tool code is enforced by
construction, at zero model-attention cost. This is what makes the harness
transfer across models: the enforced structure is identical regardless of which
model reasons over it.

\section{Develop Small, Deploy Large}

We treat model capability as a train/test axis. The small model
(\texttt{gpt-oss-120b}) is the development crucible: cheap and fast enough to run
full three-trial public-test sweeps repeatedly, so the harness can be hardened
against real failures. Two rules govern what ships.

First, a change is accepted only if a \emph{stronger} model already outperformed
a \emph{weaker} one on the bare prompt, before any domain skill. The bare-prompt
ordering (stronger $>$ weaker) is the signal that a task rewards reasoning and
structure rather than memorized scaffolding. A change that improved the small
model while regressing the large one is rejected as overfitting to the small
model's weaknesses.

Second, when the small model failed a task, we reproduced it on \texttt{gpt-5.5}
before acting. If the stronger model passed, the ceiling was reachable within the
current architecture and the fix was to encode the missing determinism as tool
code (policy as code). If both failed, the defect was in existing harness and was
repaired there rather than patched over. Candidate changes were judged on
consistent-failure sets across runs, never single noisy trials, because near the
plateau one trial moves several points on evaluator/simulator variance alone.

For Track~1 we then \emph{deploy} the identical agent --- same image, same domain
skill, no architecture change --- on \texttt{gpt-5.5}, selected purely through
environment variables. Nothing about the harness is model-specific, so the larger
model inherits the same enforced policy and normalization and applies stronger
interpretation on top.

\section{Results}

The development benchmark is the full public test split on the small model
(\texttt{gpt-oss-120b}, 125 tasks, 3 trials): \textbf{Pass\textsuperscript{1}
93.3\%} and \textbf{Pass\textsuperscript{3} 87.3\%} (per-split
Pass\textsuperscript{1}: base 90.0, hallucination 98.7, disambiguation 90.7). We
do not self-evaluate the hidden set; the numbers below characterize the deployed
\texttt{gpt-5.5} agent from our own runs. The user simulator and policy judge are
Gemini~2.5~Flash, the organizers' default.

\begin{table}[t]
\centering
\begin{tabular}{lrr}
\toprule
Per task (median) & \texttt{gpt-oss-120b} & \texttt{gpt-5.5} \\
\midrule
Model calls          & 2      & 2      \\
A2A turns            & 7      & 10     \\
Model latency / call & $\sim$1\,s & $\sim$10\,s \\
Input tokens         & 78k    & 76k    \\
Output tokens        & 1.4k   & 0.9k   \\
Thinking tokens      & \multicolumn{1}{c}{--} & 0.15k \\
\bottomrule
\end{tabular}
\caption{The call-decoupling is model-independent: two model calls carry a full
multi-turn task on both models (7 vs.\ 10 A2A turns). \texttt{gpt-5.5} figures are
measured over the OpenAI API in our runs; all three token fields, including
\texttt{thinking\_tokens}, populate.}
\label{tab:profile}
\end{table}

\paragraph{The efficiency is structural.}
Table~\ref{tab:profile} shows the coroutine bridge holds a median of two model
calls per task on \emph{both} models, even though the tasks require seven to ten
A2A tool round-trips. A conventional next-action agent would issue one model call
per round-trip. The saving is a property of the harness, not of any model.

\paragraph{Call-minimization pays off more at scale.}
The larger model's completions are roughly an order of magnitude slower
($\sim$10\,s vs.\ $\sim$1\,s per call in our measurements over the OpenAI API).
Every model call the bridge avoids therefore saves about ten seconds on
\texttt{gpt-5.5}. A per-turn baseline on the same tasks would issue roughly ten
model calls instead of two; the bridge removes on the order of eighty seconds of
sequential model latency per task. The more expensive the model, the more the
structure is worth --- the same design that trims latency on a fast small model
trims far more on a slow large one.

\paragraph{Reliability transfer.}
During development, \texttt{gpt-5.5} served as the stronger-model oracle on
residual failure clusters --- tasks curated precisely because the small model was
unreliable on them. On the final agent it resolves those targeted cases and
passes the digest-pinned GHCR validation smoke. The Open-track
Pass\textsuperscript{3} is what the organizers measure on the hidden set; our
evidence is that the harness developed on the small model deploys unchanged and
improves on the larger one.

\paragraph{Prefix-cache efficiency.}
The same one-big-prompt design that transfers policy across models also makes the
prompt cache-friendly on any provider. The system prompt is one static block
(rules, domain skill, full tool catalog); per-task state is appended only at the
tail, so the cacheable prefix is byte-identical across calls and, crucially,
\emph{across tasks} --- we never swap skills, which would break it. Across all
development runs on Cerebras (918M input tokens) 73.2\% of input tokens were
served from cache, reaching 88.7\% (p90) with the frozen prompt. OpenAI's
route applies the same automatic prefix caching, so the \texttt{gpt-5.5}
deployment inherits the same class of saving; combined with two model calls per
task, real input compute is a small fraction of the nominal per-task tokens.

\paragraph{Compute compliance.}
The internal loop is capped at five sequential model calls per step; observed
usage is a median of two (p90 four). Mean input plus output stays well under the
budget. Token and latency accounting is reported through \texttt{turn\_metrics}
on each concluding response, aggregated across intermediate tool exchanges, so
per-task totals are complete; \texttt{prompt\_tokens}, \texttt{completion\_tokens},
and \texttt{thinking\_tokens} all populate on the \texttt{gpt-5.5} route.

\section{Limitations and Conclusion}

We report a full public-test evaluation on the development model and a targeted
plus smoke validation on the deployment model; a full self-evaluation on
\texttt{gpt-5.5} was not run because the hidden set is organizer-scored and a
full sweep on a frontier model is costly. Worker state is in memory and not
recoverable across process restarts. Latency figures for \texttt{gpt-5.5} are
wall-clock over the public OpenAI API and include network and provider queueing.

The submission makes one architectural bet and reuses it across a large
capability gap: put deterministic policy in code the model calls, minimize the
number of expensive model calls, and let a small fast model do the development
that a large model then benefits from unchanged. The same image serves both the
Cerebras Track~2 and this Open-track submission, differing only by environment
variables --- direct evidence that the reliability is in the harness, not the
model.

\bibliographystyle{named}
\bibliography{track1_report}

\end{document}
